\section{Reproducibility}
\label{sec:reproducibility}

The source code, the committed documents the real-data results are read from, and the scripts of this study are in the repository named in the header block, under \texttt{manuscripts/e-process-sequential-monitors/}. The code is release v\nProductRelease\ of CAOS\_Contraste (tag v\nProductRelease, \href{https://github.com/fsantibanezleal/CAOS_Contraste/releases/tag/v\nProductRelease}{release page}). The monitors are the subpackage \texttt{riskvalidation.validation.sequential} of the engine \texttt{riskvalidation} release 1.1.0 (with its market, machine-learning and stress extras, which the other cases use), run on Python \nPythonVersion. Every number of the paper is read from a committed output by a script, and the macros of \texttt{numbers.tex} list each number with the output and the key it comes from.

Each experiment is one parameterised script run from the repository root with the interpreter of the product's environment. \texttt{experiments/exp\_hits.py} (the main VaR experiment, the calibration of the simulated boundaries and the 5,000-day size curves; \nRunHits\ seconds), \texttt{exp\_ablation.py} (\nRunAblation\ seconds), \texttt{exp\_x3\_seeds.py} (the \nSeedsCount\ seeds; \nRunSeeds\ seconds), \texttt{exp\_pd\_correlation.py} (\nRunPd\ seconds), \texttt{exp\_c11.py} and \texttt{exp\_c03.py} (the real windows and vintages; \nRunCEleven\ and \nRunCThree\ seconds),
\texttt{exp\_theory.py} (the enumerations, the exact powers and the exact spending of the monitors; \nRunTheory\ seconds), \texttt{exp\_matched.py} (the single alternative matched to each severity; \nRunMatched\ seconds) and \texttt{exp\_pd\_comparators.py} (the PD tests of the engine and the chi-square rule of equal size on the histories of the correlation and ablation experiments; \nRunPdComparators\ seconds) write \texttt{experiments/results\_*.json}. The seeds are fixed in the scripts (the base seed is \nSeed, and each scenario draws from a \texttt{SeedSequence} indexed by the seed and a scenario tag), so a rerun reproduces the committed outputs exactly; \texttt{experiments/run\_all.py} runs every script into a scratch folder, compares each output with the committed one and records the times
(\nRunTotal\ seconds in all, as wall time with other jobs running on the machine, on the hardware recorded in \texttt{run-times.json}). Then \texttt{make\_numbers.py}, \texttt{make\_tables.py} and \texttt{make\_figures.py} write \texttt{numbers.tex}, the tables and the vector figures, and \texttt{pdflatex} and \texttt{bibtex} build the document. The unit tests of the comparators, the contract tests of the outputs and the tests of the exact computations and the PD comparators are \texttt{experiments/test\_common.py}, \texttt{experiments/test\_results.py} and \texttt{experiments/test\_review\_fixes.py}.
The real-data results read the committed documents \texttt{data/derived/C11} and \texttt{data/derived/C03}, which hold exceedance days and vintage curves and not the derived-only series; the PD curves of the cohorts are recomputed from the committed vintage curves with the engine and equal the committed monitor output for \nCThreeRebuilt\ of the five cohorts.
